We ask how training-set size and the \emph{kind} of label noise jointly determine the accuracy of a small U-Net on synthetic $64\times64$ images of overlapping circles, squares and triangles, where the clean ground truth is known exactly. Training masks are corrupted either by boundary jitter or by whole-object class flips (nominal level up to 30\%); training-set sizes range from 100 to 2\,000 images; we compare plain cross-entropy (CE) with reimplemented GCE and SCE losses and a simple band-ignore loss. On clean labels mIoU grows from \rhval{analysis/analysis/flip_p0.3/ce_miou_none_n100/mean:2} at $N=100$ to \rhval{analysis/analysis/flip_p0.3/ce_miou_none_n2000/mean:2} at $N=2000$ (Welch $p=$\rhval{analysis/analysis/flip_p0.3/h1_p/mean:3}, four seeds). At $N=500$, boundary jitter at level 0.3 costs CE only \rhval{analysis/analysis/flip_p0.3/drop_boundary_n500/mean:2} mIoU, whereas object flips at the same level cost \rhval{analysis/analysis/flip_p0.3/drop_flip_n500/mean:2} (registered Welch test of flip vs boundary mIoU: $p=$\rhval{analysis/analysis/flip_p0.3/h2_p/mean:3}). More data did not dilute flip noise in our runs: the flip-induced drop is not smaller at $N=2000$ than at $N=100$ (Welch $p=$\rhval{analysis/analysis/flip_p0.3/h3_flip_p03_p/mean:2}, inconclusive), and a training-length sweep up to 1200 steps shows no late-training collapse under flips, although the network fits the flipped labels at $N=100$. No robust loss is reliably better than CE: GCE has the higher mean mIoU under flip noise (CE minus GCE: $\rhval{cmp/main/gce/flip_p0.3/miou/delta:2}$), which is not significant with four seeds ($p=$\rhval{cmp/main/gce/flip_p0.3/miou/welch_p:2}), SCE is indistinguishable from CE, and our band-ignore loss is clearly worse on clean and boundary-noise labels. This is a small CPU study (four seeds in the main grid, three in the sweeps, one architecture, one synthetic domain), and per-seed variance under flip noise is large.
